\documentclass[12pt,a4paper]{article}
\usepackage[T1]{fontenc}
\usepackage{lmodern}
\usepackage[utf8]{inputenc}
\IfFileExists{italian.ldf}{\usepackage[italian]{babel}}{}
\usepackage[margin=2.5cm]{geometry}
\usepackage{amsmath,amssymb}
\usepackage{cancel}
\usepackage{graphicx}
\usepackage{xcolor}

% Math macros (ITA)
\DeclareMathOperator{\sen}{sen}
\DeclareMathOperator{\dive}{div}
\DeclareMathOperator{\rot}{rot}
\DeclareMathOperator{\grad}{grad}
\newcommand{\defeq}{\stackrel{\mathrm{def}}{=}}

% figure riprese dalla scansione
\graphicspath{{fig/}}
\newcommand{\fig}[2]{\begin{center}\includegraphics[scale=1.2]{#1}\\[2pt] {\small fig.~#2}\end{center}}

% Ogni capitolo comincia in cima a una nuova pagina (solo \section numerate:
% \section* resta invariata, theindex la usa dentro \twocolumn[...])
\makeatletter
\let\capitolo@section\section
\renewcommand{\section}{\@ifstar{\capitolo@section*}{\clearpage\capitolo@section}}

% Indice: i titoli dei capitoli (in grassetto) senza numero di pagina,
% come nell'originale; per il resto \l@section e' quella della classe article
\renewcommand*\l@section[2]{%
  \ifnum \c@tocdepth >\z@
    \addpenalty\@secpenalty
    \addvspace{1.0em \@plus\p@}%
    \setlength\@tempdima{1.5em}%
    \begingroup
      \parindent \z@ \rightskip \@pnumwidth
      \parfillskip -\@pnumwidth
      \leavevmode \bfseries
      \advance\leftskip\@tempdima
      \hskip -\leftskip
      #1\nobreak\hfil\nobreak\hb@xt@\@pnumwidth{\hss}\par
    \endgroup
  \fi}
\makeatother

\setlength{\emergencystretch}{2em}
\setlength{\parindent}{0pt}
\setlength{\parskip}{0.4\baselineskip}
